\begin{abstract}
Der Einsatz von Reinforcement Learning zur automatisierten Zugsteuerung scheitert in der Praxis häufig an fehlender geeigneter Initialisierung: Agenten ohne fachliche Vorprägung konvergieren schlecht, während eine rein simulative Offline-Initialisierung dazu führt, dass Modelle Trainingsstrecken auswendig lernen statt eine generalisierbare Strategie zu erlernen. Diese Arbeit untersucht daher, ob sich aus historischen Betriebsdaten per Behavior Cloning ein neuronales Netz trainieren lässt, das eine generalisierbare Fahrstrategie erlernt und als Ausgangspunkt für zukünftige Reinforcement-Learning-Anwendungen dienen kann. Auf Basis realer Sensordaten einer Güterverkehrsstrecke werden eine deterministische Preprocessing-Pipeline und eine fahrtenbasierte Datenaufteilung zur Vermeidung von Data Leakage entwickelt. In einem dreistufigen Auswahlprozess wird ein LSTM-Netz (128-64-256-128, 650.369 Parameter) mit Adam-Optimierer und Smooth-L1-Verlust bestimmt und trainiert. Die Evaluation zeigt, dass das Modell eine Persistenz-Baseline bei regressionsbasierten Kennzahlen deutlich übertrifft, bei seltenen, dynamischen Manövern jedoch schwächelt, vermutlich aufgrund der begrenzten Datenbasis von rund 178 Stunden. Die Arbeit erbringt damit einen Machbarkeitsnachweis, zeigt aber zugleich, dass eine robuste Fahrstrategie eine deutlich umfangreichere Datengrundlage voraussetzt.
\end{abstract}
\selectlanguage{english}
\begin{abstract}
Applying reinforcement learning to automated train control is often hampered by the lack of suitable initialization: agents without prior domain knowledge tend to converge poorly, while purely simulation-based offline initialization risks models memorizing training routes instead of learning a generalizable strategy. This thesis therefore investigates whether a neural network can be trained on historical operational data via behavior cloning to learn a generalizable train driving strategy that can serve as a starting point for future reinforcement learning applications. Based on real sensor data from a freight rail line, a deterministic preprocessing pipeline and a trip-based data split to prevent data leakage are developed. Through a three-stage selection process, an LSTM network (128-64-256-128, 650,369 parameters) with the Adam optimizer and a Smooth-L1 loss function is determined and trained. Evaluation shows that the model clearly outperforms a persistence baseline on regression-based metrics but struggles with rare, dynamic maneuvers, likely due to the limited data basis of roughly 178 hours. The thesis thus provides a proof of concept, while also showing that a robust driving strategy requires a substantially larger data foundation.
\end{abstract}
\selectlanguage{ngerman}